\begin{blocksection}
\question Find the runtime for each of the following functions in
$\Theta$ notation. Assume \lstinline{n} is a non-negative integer and
arithmetic and printing each take constant time.

\begin{center}
\begin{tabular}{ |p{13cm}|p{2cm}| }
\hline
\begin{lstlisting}
def mystery(n):
    total = 0
    for i in range(n // 2):
        j = 1
        while j < n:
            j *= 2
            total += i * j
    return total
\end{lstlisting} &  \\  \hline

\begin{lstlisting}
def foo(n):
    if n == 0:
        return
    for i in range(n // 2):
        print(i)
    foo(n // 2)
\end{lstlisting} &  \\  \hline
\end{tabular}
\end{center}

\end{blocksection}

\begin{solution}
\begin{itemize}
    \item \lstinline{mystery}: $\Theta(n\log n)$. The outer loop runs
    $\lfloor n/2\rfloor$ times. On every iteration, \lstinline{j} resets to
    1 and doubles until it reaches or exceeds $n$, so the inner loop takes
    $\Theta(\log n)$ time. Multiply the two counts.
    \item \lstinline{foo}: $\Theta(n)$. The work across successive calls is
    proportional to $n/2 + n/4 + n/8 + \cdots$, which sums to $\Theta(n)$.
    The recursion has logarithmic depth, but its loops get shorter at each
    level; multiplying $n$ by the depth would overcount the work.
\end{itemize}
\end{solution}
